% TOP_PROBLEM: {"problemId": "problem.density-of-hyperbolicity-for-rational-maps", "canonicalTitle": "Density of Hyperbolicity for Rational Maps", "releaseRank": 155, "primaryDomain": "probability_ergodic_dynamics", "primaryDomainLabel": "Probability, ergodic theory & dynamics", "displayStatus": "open", "releaseStatus": "open"}
% !TeX program = lualatex
% Definition and sources only; this file is not an LLM solution attempt.
\documentclass[11pt]{article}
\usepackage[margin=25mm]{geometry}
\usepackage{fontspec}
\setmainfont{Latin Modern Roman}
\usepackage{amsmath,amssymb,mathtools}
\usepackage{unicode-math}
\setmathfont{Latin Modern Math}
\usepackage{xurl,hyperref}
\hypersetup{colorlinks=true,urlcolor=blue}
\setcounter{secnumdepth}{0}
\setlength{\parindent}{0pt}
\setlength{\parskip}{5pt}
\setlength{\emergencystretch}{3em}
\begin{document}
\section*{\texorpdfstring{Density of Hyperbolicity for Rational Maps}{Density of Hyperbolicity for Rational Maps}}\label{155}
\subsection{Definitions and mathematical statement}
Represent a degree-$d$ rational map by $f(z)=P(z)/Q(z)$ with coprime polynomials, modulo common nonzero scalar multiplication, giving the parameter space $\operatorname{Rat}_d\subset{\mathbb{P}}^{2d+1}({\mathbb{C}})$. A critical point is a point of the sphere where the local mapping degree exceeds one. A periodic cycle is attracting if the modulus of its return-map multiplier is less than one. Call $f$ hyperbolic when every critical orbit converges to an attracting cycle. For each $d\ge2$, ask whether
\[
 \overline{\{f\in\operatorname{Rat}_d:f\text{ is hyperbolic}\}}=\operatorname{Rat}_d
\]
in the usual complex parameter topology. The target is density, not that every map is hyperbolic or structurally stable.
\par\smallskip\textit{Formulation and concept references:} \href{https://arxiv.org/abs/1709.09869}{[S1]}.

\subsection{Short English statement}
Can every rational dynamical system of a fixed degree be approximated arbitrarily closely by one whose critical points all settle into attracting cycles?
\subsection{Sources}
\begin{itemize}
\item[S1]A. M. Benini, 'A survey on MLC, Rigidity and related topics', arXiv:1709.09869 (2017, revised 2018).
\url{https://arxiv.org/abs/1709.09869}.
\textit{Formulation source.}
\end{itemize}
\end{document}
